\scanpage{21}
\uline{Re N.B.} Soit tjrs $g=$ dim rel. de $A/S$, et considérons
$\omega_A=\deleted{A}\operatorname{d\acute et}(\underline t_A)^{-1}
=\bigwedge^g\underline t_A\deleted{\unsure{{}^\vee}}$
\inserted{(formes diff de degré max le long de la section unité)}.
C'est un faisceau inv.
à p-structure non dégénérée, i.e. déduit par
ext. des scalaires d'un \deleted{\unsure{torseur}}\deleted{\unsure{sur}}\inserted{\deleted{\unsure{fais|faisc.}} faisceau loc. ct}
\deleted{\unsure{$\mathbb F_p^*$|$({}_pA^*)_{\mathrm{\acute et}}$}}\unsure{de} \unsure{inv|dim} sur $\mathbb F_p$, savoir
$\deleted{\unsure{(|\{}}\operatorname{d\acute et}\unsure{{}_g|{}_p|{}}({}_pA^*)_{\mathrm{\acute et}}\deleted{\unsure{)|\}}}$.
À ce titre, $\omega_A^{\otimes(p-1)}$ admet une section canonique.

Si on ne suppose plus \unsure{néc.} $A/S$ ordinaire, et si $U$ est l'ouvert des $s\in S$
tels que $A_s$ soit ordinaire, tt ce qui
précède s'applique à $A|U$, en particulier
on trouve une section canonique \inserted{$\sigma$} de
$L|U$, partout $\ne0$. \deleted{Si $U$ est sch. dense dans $S$,
on peut considérer $\sigma$ comme une \deleted{\unsure{fon}}\inserted{pseudo section}
\unsure{de $L$} sur $S$ (ou fonction méromorphe sur $S$), \deleted{\unsure{qui se}}\inserted{mais en}}

Mais en fait, \unsure{elle|on} provient d'une section
canonique de $L^{\otimes(p-1)}$ sur $S$ tt entier, défini
à l'aide de
\[
t_A^{(p)}\xrightarrow{\pi}t_A
\]
en passant au dét, d'où
\[
(\operatorname{d\acute et}\underline t_A)^{\otimes p}\longrightarrow\operatorname{d\acute et}\underline t_A
\]
d'où
\[
\sigma_A\in\Gamma(\operatorname{det}\underline t_A)^{1-p}=\Gamma\omega_A^{(p-1)}
\]
% archive sheet number: 20
